%% Работа с русским языком
\usepackage{cmap}			 % поиск в PDF
\usepackage{mathtext} 		 % русские буквы в формулах
\usepackage[T2A]{fontenc}	 % кодировка
\usepackage[utf8]{inputenc}	 % кодировка исходного текста
\usepackage[russian]{babel}	 % локализация и переносы
\usepackage[htt]{hyphenat}
\usepackage{microtype}
\emergencystretch=3em

%% Пакеты для работы с математикой
\usepackage{amsmath,amsfonts,amssymb,amsthm,mathtools}
\usepackage{icomma}

%% Нумерация формул (опционально)
%\mathtoolsset{showonlyrefs=true} % показывать номера только у тех формул, на которые есть \eqref{} в тексте.
%\usepackage{leqno}               % нумерация формул слева

%% Шрифты
\usepackage{euscript}	 % шрифт "Евклид"
\usepackage{mathrsfs}    % красивый мат. шрифт

%% Некоторые полезные макросы для дебага (в случае недоверия авторам шаблона)
\makeatletter
\newcommand\thefontsize{The current font size is: \f@size pt} % пример: \section{\thefontsize}
\makeatother

%% Настройка размеров шрифтов
\makeatletter
\setlength{\headheight}{28pt}
%% TODO: мне не удалось разобраться, как грамотно подбирать второе число в
%% \@setfontsize\*, но ряд экспериментов показывает, что "10" выравнивает текст весьма прилично :)
\renewcommand\Huge{\@setfontsize\Huge{14pt}{10}}
\renewcommand\huge{\@setfontsize\huge{14pt}{10}}
\renewcommand\Large{\@setfontsize\Large{14pt}{10}}
\renewcommand\large{\@setfontsize\large{12pt}{10}}
\makeatother

%% Поля (геометрия страницы)
\usepackage[left=3cm,right=1.5cm,top=2cm,bottom=2cm,bindingoffset=0cm]{geometry}

%% Русские списки
\usepackage{enumitem}
\makeatletter
\AddEnumerateCounter{\asbuk}{\russian@alph}{щ}
\makeatother

%% Работа с картинками
\usepackage{caption} % пакет для подписей к рисункам, в частности, для работы caption*
\captionsetup{justification=centering} % центрирование подписей к картинкам
\usepackage{svg}
\usepackage{graphicx}                  % вставки рисунков
\graphicspath{{images/}{images2/}}     % папки с картинками
\setlength\fboxsep{3pt}                % отступ рамки \fbox{} от рисунка
\setlength\fboxrule{1pt}               % толщина линий рамки \fbox{}
\usepackage{wrapfig}                   % обтекание рисунков и таблиц текстом

%% Работа с таблицами
\usepackage{array,tabularx,tabulary,booktabs} % дополнительная работа с таблицами
\usepackage{longtable}                        % длинные таблицы
\usepackage{multirow}                         % слияние строк и столбцов в таблице
\usepackage{makecell}                         % создание ячеек с многими строчками текста
\usepackage{hhline}                           % позволяет разрывать горизонтальные линии в таблицах

%% Красная строка
\setlength{\parindent}{1.25cm}

%% Интервалы
\linespread{1}

%% TikZ
\usepackage{tikz}
\usetikzlibrary{graphs,graphs.standard}

%% Верхний колонтитул
\usepackage{fancyhdr}
\pagestyle{fancy}

%% Перенос знаков в формулах (по Львовскому)
\newcommand*{\hm}[1]{#1\nobreak\discretionary{}{\hbox{$\mathsurround=0pt #1$}}{}}

%% Дополнительно
\usepackage{float}   % добавляет возможность работы с командой [H] которая улучшает расположение на странице
\usepackage{gensymb} % красивые градусы
\usepackage{listings} % пакет для листингов с кодом

\lstset{
    language=C,
    % рамка вокруг листинга
    frame=single,
    % номера строк слева
    numbers=left,
    % расстояние между номерами строк и кодом
    numbersep=8pt,
    % подпись под листингом
    captionpos=b,
    % таб в коде равен 4 пробелам
    tabsize=4,
    % делает кавычки в моноширинном шрифте нормальными, полезно
    upquote=true,
    % перенос слишком длинных строк
    breaklines=false,
    % показывать табы в коде
    showtabs=false,
    % показывать пробелы в коде
    showspaces=false,
    % показывать пробелы в строковых литералах в коде
    showstringspaces=false,
    % настройки стиля, применяемые ко всему коду
    basicstyle=\fontsize{10}{9.1}\selectfont\ttfamily,
    % настройки стиля ключевых слов
    keywordstyle=\color{blue}\bfseries,
    % настройки стиля строковых литералов
    stringstyle=\color{red},
    % настройки стиля номеров строк
    numberstyle=\tiny\color{gray}
}

\definecolor{CommentGreen}{RGB}{0,150,0}
\lstdefinestyle{normal}{
    % настройки стиля комментариев
    commentstyle=\color{CommentGreen},
}

\definecolor{DiffGreen}{RGB}{0,150,0}
\definecolor{DiffRed}{RGB}{200,0,0}
\lstdefinestyle{diff}{
    commentstyle=,
    moredelim=[is][\color{DiffGreen}]{<+>}{</+>},
    moredelim=[is][\color{DiffRed}]{<->}{</->}
}

\definecolor{DeltaGreen}{RGB}{0,150,0}
\definecolor{DeltaRed}{RGB}{200,0,0}
\definecolor{DeltaMid}{RGB}{184,134,11}

% Hyperref (для ссылок внутри  pdf)
\usepackage[unicode, pdftex]{hyperref}

% Отступ перед первым абзацем в каждом разделе
\usepackage{indentfirst}
